%\documentclass{article}
%\usepackage[utf8]{inputenc}
%\usepackage[english]{babel}
%\usepackage[]{amsthm} %lets us use \begin{proof}
%\usepackage[]{amssymb} %gives us the character \varnothing
%\usepackage{verbatim}

\documentclass[11pt]{article}
\usepackage{amsmath,amssymb,amsthm,latexsym}
\usepackage{mathrsfs,mathtools}
\usepackage{graphicx}
\usepackage{comment}
\usepackage{bm}
\usepackage{natbib}
\usepackage{subfigure,multirow}
\usepackage[colorlinks,linkcolor=black,citecolor=black,urlcolor=black]{hyperref}
\usepackage{verbatim}

\renewcommand\baselinestretch{1.25}
\topmargin -.5in
\textheight 9in
\textwidth 6.6in
\oddsidemargin -.25in
\evensidemargin -.25in

\title{Solutions to Homework 2}
\author{Xiaochuan Gong}
\date{March 11, 2022}
%This information doesn't actually show up on your document unless you use the maketitle command below

\begin{document}
\maketitle %This command prints the title based on information entered above

\subsection*{Problem 1}
For a function $f: {\mathbb R}^n \rightarrow {\mathbb R}\cup\{-\infty, +\infty\}$ the following statements are equivalent:

\noindent (a) $f$ is closed

\noindent (b) $C_\alpha$ is closed for any $\alpha$

\noindent (c) $f$ is lower-semicontinous over ${\mathbb R}^n$
\begin{proof}
(a)$\Rightarrow$(b) Let $\alpha$ be any scalar and consider $C_\alpha$. If $C_\alpha=\emptyset$, then $C_\alpha$ is closed. Suppose now that $C_\alpha\neq\emptyset$. Pick $\{x_k\}\subset C_\alpha$ such that $x_k \rightarrow \bar x$ for some $\bar x \in {\mathbb R}^n$. We have $f(x_k)\leq\alpha$ for all $\alpha$, implying that $(x_k, \alpha)\in \mathrm{epi}(f)$ for all $\alpha$. Since $(x_k, \alpha) \rightarrow (\bar x, \alpha)$ and $\mathrm{epi}(f)$ is closed, it follows that $(\bar x, \alpha) \in \mathrm{epi}(f)$. Consequently $f(\bar x)\leq\alpha$, showing that $\bar x \in C_\alpha$. Hence $C_\alpha$ is closed for any $\alpha$.\\

\noindent (b)$\Rightarrow$(c) Let $x_0 \in {\mathbb R}^n$ be arbitary and let $\{x_k\}$ ba a sequence such that $x_k \rightarrow x_0$. To arrive at a contradiction, assume that $f$ is not lower-semicontinous at $x_0$, i.e.,
$$\mathop{\lim\inf}\limits_{k \rightarrow \infty} f(x_k) < f(x_0).$$
Then, there exist a scalar $\beta$ and a subsequence $\{x_{k_j}\}\subset \{x_k\}$ such that
$$f(x_{k_j})\leq\beta < f(x_0),$$
yielding that $\{x_{k_j}\} \subset C_\beta$. Since $x_k \rightarrow x_0$, then $x_{k_j}\rightarrow x_0$. Note that the set $C_\beta$ is closed, it follows that $x_0 \in C_\beta$. Hence, $f(x_0)\leq\beta$, a controdiction arrives. Thus, we must have
$$f(x_0)\leq\mathop{\lim\inf}\limits_{k \rightarrow \infty} f(x_k).$$
\\
\noindent (c)$\Rightarrow$(a) To arrive at a controdiction we assume that $\mathrm{epi}(f)$ is not closed. Then, there exists a sequence $\{(x_k, t_k)\}\subset \mathrm{epi}(f)$ such that 
$$(x_k, t_k) \rightarrow (\bar x, \bar t) \quad\text{and}\quad (\bar x, \bar t) \notin \mathrm{epi}(f).$$
Since $(x_k, t_k) \in \mathrm{epi}(f)$ for all $k$, we have
$$f(x_k) \leq t_k, \forall k.$$
Taking the limit inferior as $k \rightarrow \infty$, and using $t_k \rightarrow \bar t$, we obtain
$$\mathop{\lim\inf}\limits_{k \rightarrow \infty} f(x_k) \leq \lim\limits_{k \rightarrow \infty} t_k = \bar t.$$
Since $(\bar x, \bar t) \notin \mathrm{epi}(f)$, we have $f(\bar x) > \bar t$, implying that
$$\mathop{\lim\inf}\limits_{k \rightarrow \infty} f(x_k) \leq \bar t < f(\bar x).$$
On the other hand, because $x_k \rightarrow \bar x$, and $f$ is lower-semicontinous at $\bar x$, we have
$$\mathop{\lim\inf}\limits_{k \rightarrow \infty} f(x_k)\geq f(\bar x),$$
a controdiction arrives. Hence, $\mathrm{epi}(f)$ must be closed.

\end{proof}

\end{document}